\documentclass[10pt,a4paper]{amsart}
\usepackage[margin=19mm]{geometry}
\usepackage{amsmath,amssymb,mathtools}
\usepackage[scr=rsfs,cal=euler]{mathalfa}
\usepackage{xfrac}
\usepackage[unicode,hidelinks,bookmarks=false]{hyperref}
\newcommand{\ie}{i.e.\ }
\newcommand{\cf}{cf.\ }
\newcommand{\resp}{respectively\ }
\newcommand{\Subsection}{\subsection}
\providecommand{\nobreakdash}{\nobreak\hbox{-}}
\setlength{\emergencystretch}{2em}
\multlinegap0pt
\usepackage[shortlabels]{enumitem}
\usepackage{tikz}
\newtheorem{theorem}{Theorem}
\newtheorem{proposition}{Proposition}
\providecommand{\F}{\mathbb{F}}
\title{On Rado's single equation theorem}
\author{Tom Sanders}
\date{Source: arXiv:2603.18179v1 (CC BY 4.0)}
\begin{document}
\maketitle

\noindent\emph{Source and licence.} On Rado's single equation theorem. Version arXiv:2603.18179v1 (CC BY 4.0), distributed by arXiv under the Creative Commons Attribution 4.0 International licence, http://creativecommons.org/licenses/by/4.0/, as recorded in the arXiv licence metadata for this exact version and shown on the abstract page https://arxiv.org/abs/2603.18179v1. Full licence notice: \texttt{licenses/arxiv-2603.18179-CC-BY-4.0.txt}.\par\smallskip

\noindent\textit{Editorial scope.} Counted unit: the article's theorem thm.model (source0.tex lines 132--134). The article's complete original proof of it is delivered with it (source0.tex lines 154--201, one \texttt{proof} environment of 48 lines, which the article places away from the statement). No mathematics is written, repaired or re-derived: the statement and the proof are the article's own lines, in the article's own notation, and no retained line is substituted or re-flowed.  Retained uncounted as the source-local context the proof needs: lines 146--151.  Nothing was omitted from any retained span, so the package carries no omission record.  No retained line contains a citation, so the wrapper carries no bibliography and no bibliography entry.  No retained line contains a figure environment, an image-inclusion command or drawing code, so no image is delivered and the package declares no asset dependency.  Editorial formatting only, in addition: an AMS \texttt{amsart} preamble that restates the article's own declarations and macros used by the retained lines, namely the environments \texttt{theorem}, \texttt{proposition} on separate counters, together with the standard packages those lines need.  The retained environments print in this document as Theorem 1, Proposition 1 in order of appearance, whereas the article prints the same items with the numbering of its own full document.  The complete native source of the article as distributed in the arXiv e-print for this exact version is bundled unchanged as \texttt{sources/196673/source0.tex}.
\begin{proposition}\label{prop.sum} Suppose that $V \leq G$, $A,D \subset V$ are non-empty with $\alpha:=\mu_V(A)$ and $\delta:=\mu_V(D)$, and $\epsilon \in (0,1]$ is a parameter. Then either
\begin{equation*}
|\langle 1_{-A}  \ast  1_{A}, 1_D\rangle_{L_2(\mu_G)}-\alpha^2\delta\mu_G(V)^2| \leq \epsilon \alpha^2\delta\mu_G(V)^2; 
\end{equation*}
 or else there is $V' \leq V$ with $\dim V/V' = O_\epsilon((\log^4 2\alpha^{-1})(\log^4 2\delta^{-1}))$ such that $\|1_A \ast \mu_{V'}\|_\infty \geq (1+\Omega(\epsilon))\alpha$.
\end{proposition}

\begin{theorem}\label{thm.model}
Suppose that $a,b \in \F_q^*$ and $\mathcal{C}$ is an $r$-colouring of $G$. Then there is a colour class $A \in \mathcal{C}$ and $\exp(-O_q(r\log^92r)) |G|^2$ triples $(x,y,z) \in A^3$ such that $ax-ay = bz$.
\end{theorem}

\begin{proof}[Proof of Theorem \ref{thm.model}]
Write $\mathcal{C}:=\{A_1,\dots,A_r\}$. We proceed iteratively to construct $G=:V_0 \geq V_1\geq \dots$. Writing $\cdot$ for the scalar multiplication of $\F_q$ on $G$, set
\begin{equation*}
S_{i,j}:=\|1_{a\cdot A_j}\ast \mu_{V_i}\|_\infty \text{ for all } j \in [r].
\end{equation*}
Since $V_i \leq V_{k}$ whenever $i \geq k$ we have $\mu_{V_i} \ast \mu_{V_k} = \mu_{V_k}$ and hence
\begin{equation}\label{mon}
S_{i,j} \geq S_{k,j}\text{ whenever }i \geq k, \text{ and } S_{i,j} \leq 1 \text{ for all }i.
\end{equation}

At step $i$ of the iteration we shall show that either
\begin{enumerate}
\item\label{c1} there is some $j \in [r]$ such that
\begin{equation*}
\langle 1_{a\cdot A_j} \ast 1_{-a\cdot A_j},1_{b \cdot A_j}\rangle_{L_2(\mu_G)} \geq \frac{1}{2r^3}\mu_G(V_i)^2;\text{ or }
\end{equation*}
\item\label{c2} there is $V_{i+1} \leq V_i$ and $j \in [r]$ such that
\begin{enumerate}
\item\label{c21} $\dim V_i/V_{i+1} =O(\log^82r)$;
\item\label{c22} $S_{i+1,j} \geq (1+\Omega(1))S_{i,j}$; and
\item\label{c23} $S_{i,j} \geq 1/r$.
\end{enumerate}
\end{enumerate}
We stop the iteration the first time we are in case (\ref{c1}). Suppose that we are in case (\ref{c2}) for a particular $j$ at steps $i_1 <i_2<\dots < i_k$ of the iteration. Then by (\ref{mon}) and (\ref{c22}) \& (\ref{c23}) we have
\begin{align*}
1 \geq S_{i_k+1,j} \geq (1+\Omega(1))S_{i_k,j} & \geq (1+\Omega(1))S_{i_{k-1}+1,j}\\ & \geq \qquad \dots \qquad  \geq (1+\Omega(1))^kS_{i_1,j} \geq (1+\Omega(1))^k\cdot (1/r).
\end{align*}
It follows that $k=O(\log 2r)$. Since there are $r$ possible values for $j$ we conclude that we must have terminated the iteration at step $i_0=r\cdot O(\log 2r)$, in which case
\begin{equation*}
\dim G/V_{i_0} =\dim V_0/V_1 + \cdots + \dim V_{i_0-1}/V_{i_0} =i_0\cdot O(\log^82r) = O(r\log^92r)
\end{equation*} by (\ref{c21}).

Since $\mu_G(V_{i_0})=q^{-\dim G/V_{i_0}}$ and $|G|^2\langle 1_{a\cdot A_j} \ast 1_{-a\cdot A_j},1_{b \cdot A_j}\rangle_{L_2(\mu_G)} $ is exactly the number of solutions to $ax-ay=bz$ for $x,y,z \in A_j$ we have the conclusion of Theorem \ref{thm.model} from (\ref{c1}) as claimed.

It remains to show that at each stage of the iteration we are either in case (\ref{c1}) or (\ref{c2}). Suppose we are at step $i$ of the iteration. Since $b \in \F_q^*$ we have $b\cdot V_i=V_i$ and $\mathcal{C}$ is a cover of $G$, and there is some $j=j(i) \in [r]$ such that $\mu_{V_i}(b\cdot A_j) \geq \frac{1}{r}$. Since $a,b \in \F_q^*$ we have $(ab^{-1})V_i=V_i$, and
\begin{equation*}
S_{i,j}=\|1_{a\cdot A_j}\ast \mu_{V_i}\|_\infty \geq \mu_{V_i}(a\cdot A_j) = \mu_{V_i}(b \cdot A_j).
\end{equation*}
In particular $S_{i,j} \geq 1/r$. Let $x_i\in G$ be such that 
\begin{equation*}
\mu_{V_i}(x_i-a \cdot A_j )=1_{a \cdot A_j} \ast \mu_{V_i}(x_i)=\|1_{a\cdot A_j}\ast \mu_{V_i}\|_\infty.
\end{equation*}
Apply Proposition \ref{prop.sum} to $V_i\leq G$ with $\epsilon:=\frac{1}{2}$, and $A:=(x_i-a \cdot A_j)\cap V_i$ and $D:=(b \cdot A_j)\cap V_i$ to get that either
\begin{align*}
\langle 1_{a \cdot A_j} \ast 1_{-a\cdot A_j}, 1_{b \cdot A_j}\rangle_{L_2(\mu_G)} & =  \langle 1_{-(x_i-a \cdot A_j)} \ast 1_{x_i-a\cdot A_j}, 1_{b \cdot A_j}\rangle_{L_2(\mu_G)}\\ &\geq \langle 1_{-A} \ast 1_{A},1_D\rangle_{L_2(\mu_{G})}\\ &\geq \frac{1}{2}\mu_{V_i}(A)^2\mu_{V_i}(D)\mu_G(V_i)^2 \geq \frac{1}{2r^3}\mu_G(V_i)^2;
\end{align*}
or else there is $V_{i+1} \leq V_i$ with $\dim V_i/V_{i+1}=O(\log^82r)$ and $S_{i+1,j} \geq (1+\Omega(1))S_{i,j}$.  In the first case we are in (\ref{c1}); in the second we are in (\ref{c2}). The result is proved.
\end{proof}

\end{document}
